%! Comparing Studies and Choosing a Method — Unit 1, Lesson 1.7 — Warm-Up (Answer Key)
% ONE PAGE, blank and key. Three items at 12pt (the stats-model size), \workrowsep 50pt, item
% spacing 22pt, so each item gets real handwriting room and still fits. Do not add a fourth
% item — the page is full.
% GRADUAL RELEASE: the warm-up activates prior knowledge before the guided notes. Items 1 and 2
% are Lesson 1.5's Riverbend arithmetic — the same 39-of-60 and 72-of-90 the notes' row 1
% display carries, so the two percents are already on the page when the comparison starts.
% Item 3 puts the two studies side by side WITHOUT the vocabulary: the two gaps, and the tell
% (60/90 against 60/60) — row 1 names the designs and problem 4 asks which result to believe.
\documentclass[12pt]{article}
\usepackage{saar-article}
\usepackage{saar-key}   % pulls in -boxes; do NOT also load -boxes
\newcommand{\TallMath}[1]{$\displaystyle #1\rule[-1.4em]{0pt}{3.2em}$}
% Word bank strip for every fill-in cluster (user request 2026-09-06). Defined per document;
% the key inherits it and prints the same strip, so heights match.
\newcommand{\wordbank}[1]{\par\vspace{2pt}\noindent\fcolorbox{goldacc}{goldbg}{\parbox{\dimexpr\linewidth-2\fboxsep-2\fboxrule\relax}{\small\textbf{Word bank:} #1}}\par\vspace{4pt}}
% Handwriting room inside every \begin{work} block. Moves the blank and the
% key together, so the two can never drift apart.
\setlength{\workrowsep}{50pt}

\begin{document}

\pageheader{Unit 1, Lesson 1.7}{Warm-Up --- Answer Key}
% NAMESTRIP: no \namedateperiod here — mirrors the blank. NO teachernote here —
% teacher prose lives in the lesson plan.

\vspace{0.15in}

\begin{notesbox}{Spiral Review --- Two Studies, One Study Center}
\normalsize
Riverbend High School has studied its after-school Study Center twice.
\textbf{Last quarter the counselor watched:} she pulled $150$ student records --- $60$
students had used the Study Center and $90$ had not. $39$ of the users passed every class;
$72$ of the non-users did. \textbf{This quarter she assigned:} a random number generator
split $120$ volunteers into two groups of $60$; $75\%$ of the treated group and $55\%$ of the
control group passed every class.

\begin{enumerate}[leftmargin=1.8em, itemsep=22pt, topsep=10pt]

  \item Last quarter, what \textbf{percent} of the $60$ Study Center users passed every class?

        \begin{work}
          \dfrac{39}{60} &= 0.65 \\
                         &= 65\%
        \end{work}

  \item Last quarter, what \textbf{percent} of the $90$ non-users passed every class?

        \begin{work}
          \dfrac{72}{90} &= 0.80 = 80\%
        \end{work}

  \item Put the two studies side by side.
        \wordbank{15 \;\textbullet\; 20 \;\textbullet\; 60 and 90 \;\textbullet\; 60 and 60 \;\textbullet\; the students themselves \;\textbullet\; a random number generator}

        \vspace{4pt}
        \renewcommand{\arraystretch}{2.6}
        \begin{tabularx}{\linewidth}{@{} X p{2.6cm} p{2.6cm} @{}}
        \toprule
        \textbf{The study} & \textbf{Gap (points)} & \textbf{Group sizes} \\
        \midrule
        Last quarter --- watched $150$ records  & \ans{15} & \ans{60 and 90} \\
        This quarter --- assigned $120$ volunteers & \ans{20} & \ans{60 and 60} \\
        \bottomrule
        \end{tabularx}

        \vspace{10pt}
        Who put each student in a group \emph{last} quarter? \ans{the students themselves}

\end{enumerate}
\end{notesbox}

\end{document}
